% !TeX program = lualatex
% !TeX encoding = utf8
% !TeX spellcheck = uk_UA

%% --------------------------------------------------------
\section{Основні позначення}\label{sec:Notations}\hypertarget{Notations}{}
%% --------------------------------------------------------

Векторні величини позначено напівжирними літерами, їхні модулі~--- звичайними. Формули записано в~гауссовій системі
одиниць; частину~VII (електричні кола, довгі лінії, хвилеводи) викладено в~СІ (див. передмову). Деякі літери в~різних
розділах означають різні величини~--- це вказано в~останньому стовпчику.

\begin{longtblr}[
		caption = {Основні позначення},
		label = {tab:Notations},
	]{
		colspec={Q[c,m,2.2cm] X[l,m] X[l,m]},
		rowhead = 1,
		row{1} = {c, fg=white, bg=themecolorlight, font=\bfseries\small},
		row{even[2]} = {gray!8},
		cells = {font=\small},
		hlines = {gray!40},
	}
	Позначення & Величина & Примітка \\
	\SetCell[c=3]{c, bg=gray!20} \textit{Електричне поле} & & \\
	$q$, $Q$ & заряд & \\
	$\rho$, $\sigma$ & об'ємна й~поверхнева густина заряду & $\rho$ також густина речовини \\
	$\lambda$ & лінійна густина заряду & у~розділах про струм $\lambda$~--- питома провідність \\
	$\Efield$, $\phi$ & напруженість і~потенціал електричного поля & \\
	$\Dfield$ & електрична індукція (зміщення) & $\Dfield=\Efield+4\pi\vect P$ \\
	$\vect P$ & вектор поляризації (поляризованість) & \\
	$\vect p$ & дипольний момент & \\
	$\alpha$ & поляризовність атома чи молекули & \\
	$\chi$ & діелектрична сприйнятливість & $\vect P=\chi\Efield$, $\epsilon=1+4\pi\chi$ \\
	$\epsilon$ & діелектрична проникність & \\
	$C$ & ємність & \\
	$W$, $w$ & енергія та її густина & \\
	\SetCell[c=3]{c, bg=gray!20} \textit{Постійний струм} & & \\
	$I$, $\vect j$ & сила й~густина струму & \\
	$\lambda$ & питома провідність & $\vect j=\lambda\Efield$; у~розділі~\ref{ElMagWaves}~--- $\sigma$ \\
	$R$, $\emf$ & опір, електрорушійна сила (ЕРС) & \\
	$\varrho$ & питомий опір & $\varrho=1/\lambda$ \\
	$\tau_\text{М}$ & максвелівський час релаксації & $\tau_\text{М}=\epsilon/4\pi\lambda$ \\
	\SetCell[c=3]{c, bg=gray!20} \textit{Магнітне поле} & & \\
	$\Bfield$, $\Hfield$ & магнітна індукція і~напруженість магнітного поля & $\Bfield=\Hfield+4\pi\vect J$ \\
	$\vect A$ & векторний потенціал & $\Bfield=\rot\vect A$ \\
	$\vect J$ & намагніченість & \\
	$\vect p_m$ & магнітний момент & \\
	$\chi_m$ & магнітна сприйнятливість & $\vect J=\chi_m\Hfield$, $\mu=1+4\pi\chi_m$ \\
	$\mu$ & магнітна проникність & \\
	$\Phi$ & магнітний потік & \\
	$L$, $L_{12}$ & індуктивність, взаємна індуктивність & $\Phi=LI/c$; у~частині~VII (СІ) $\Phi=LI$ \\
	\SetCell[c=3]{c, bg=gray!20} \textit{Хвилі, випромінювання, відносність} & & \\
	$\vect\Pi$ & вектор Пойнтінга & $\vect\Pi=\frac{c}{4\pi}\Efield\times\Hfield$ \\
	$\vect g$ & густина імпульсу поля & $\vect g=\vect\Pi/c^2$ \\
	$\vect k$, $\omega$, $\lambda$ & хвильовий вектор, колова частота, довжина хвилі & \\
	$n$ & показник заломлення & $n=\sqrt{\epsilon\mu}$ \\
	$\vect\beta$, $\gamma$ & $\vect v/c$ і~лоренців множник & $\gamma=1/\sqrt{1-\beta^2}$ \\
	\SetCell[c=3]{c, bg=gray!20} \textit{Коливання, кола, довгі лінії (частина VII, СІ)} & & \\
	$\gamma$ & коефіцієнт згасання & $\gamma=R/2L$; не плутати з~лоренцевим $\gamma$ \\
	$\omega_0$, $\delta$, $Q$ & власна частота, логарифмічний декремент, добротність & $Q$ використано лише для добротності \\
	$Z$ & імпеданс & $Z=R+i(\omega L-1/\omega C)$ \\
	$\overline P$, $Q_\text{р}$ & середня й~реактивна потужність & \\
	$R_0$, $L_0$, $C_0$, $G_0$ & погонні параметри довгої лінії & на одиницю довжини \\
	$\gamma$, $Z_\text{хв}$, $\Gamma$ & стала поширення, хвильовий опір, коефіцієнт відбиття & \\
\end{longtblr}
